\chapter{Theoretical Background}
\label{sec:chapterTheoreticalBackground}

The State of the Art chapter (Chapter~\ref{sec:stateoftheart}) surveyed the
\emph{literature landscape} around fraud detection under extreme class
imbalance: which datasets are standard, which evaluation pitfalls recur, which
debates remain open. The Methodology chapter (Chapter~\ref{sec:chapterMethodology})
will describe how the present study was actually conducted: which datasets,
which splits, which hyperparameter budget, which leakage-free protocol. Sitting
between the two, this chapter introduces the \emph{theoretical toolkit} that
the methodology draws on. Each tool --- learning algorithm, imbalance
intervention, evaluation metric, threshold rule, optimiser, attribution
estimator --- is defined here once, in general terms, so that the Methodology
chapter can focus exclusively on project-specific design choices and refer
back to this chapter whenever a definition or mechanism needs to be recalled.

The exposition deliberately stops short of comparative analysis. Which
combinations perform best on the present datasets is an empirical question
answered in the Results and Discussion chapter
(Chapter~\ref{sec:chapterResults}). The narrower purpose here is to provide a
self-contained reference for the terminology and mechanisms used in the
remainder of the thesis.

% ──────────────────────────────────────────────────────────────────────────
\section{Model Families}
\label{sec:tb:models}
% ──────────────────────────────────────────────────────────────────────────

The study covers four supervised approaches --- a linear classifier, bagged
trees, gradient-boosted trees, and an attention-based deep network --- together
with a one-class anomaly-detection baseline. This section recalls their defining
mechanisms; the rationale for selecting these families is given in the
Methodology chapter.

\paragraph{Logistic Regression.}
Logistic Regression is the canonical linear probabilistic classifier. It
models the log-odds of the positive class as a linear combination of the input
features, $\log\!\frac{P(y{=}1\,|\,\mathbf{x})}{P(y{=}0\,|\,\mathbf{x})} = \mathbf{w}^{\top}\mathbf{x} + b$,
and recovers a probability through the sigmoid link $\sigma(z)=(1+e^{-z})^{-1}$.
Parameters are typically learned by minimising the regularised
cross-entropy loss; the elastic-net penalty
$\lambda(\alpha\|\mathbf{w}\|_1 + (1-\alpha)\|\mathbf{w}\|_2^2/2)$ combines
$\ell_1$ and $\ell_2$ regularisation and acts as a unified sparsity/shrinkage
control. Logistic Regression is interpretable through its coefficients, can be
well-calibrated when the modelling assumptions are adequate, and serves as a
simple parametric reference against which more flexible models are compared.

\paragraph{Random Forest.}
A Random Forest is a bagged ensemble of decision trees. Each tree is fitted on a
bootstrap sample of the training data and, at every split, considers only a
random subset of features; the ensemble prediction is the average of the
per-tree class probabilities. The combination of bootstrap aggregation and
feature subsampling is designed to reduce correlation among trees and lower
ensemble variance relative to a single tree. Tree splits are also relatively
insensitive to monotonic feature scaling.

\paragraph{Gradient-boosted decision trees.}
Gradient boosting builds an additive ensemble
$F_{m}(\mathbf{x}) = F_{m-1}(\mathbf{x}) + \nu\, h_m(\mathbf{x})$ in which
each new weak learner $h_m$ is fit to the negative gradient of a chosen loss
with respect to the current prediction, scaled by a learning rate $\nu$. Two
modern implementations are used in this work. \textbf{LightGBM}~\cite{ke2017lightgbm}
adopts a \emph{leaf-wise} (best-first) tree growth strategy, which splits the
single leaf with the largest loss reduction at each step rather than expanding
the tree level by level, and bins continuous features into discrete histograms
so that split evaluation becomes a constant-time operation per bin.
\textbf{CatBoost}~\cite{prokhorenkova2018catboost} introduces \emph{ordered
boosting}, which computes the residual for each training example using a model
constructed without that example, thereby reducing the prediction-shift bias
that can arise when the same data are used to compute gradients and fit the
next learner. CatBoost also handles categorical
features natively via permutation-aware target statistics, removing the need
for one-hot encoding of high-cardinality categorical inputs.

\paragraph{One-Class SVM.}
The One-Class Support Vector Machine is an unsupervised anomaly-detection
model. Given training examples drawn predominantly from a single class, it
maps the observations into a kernel feature space and learns a maximum-margin
hyperplane that separates most training observations from the origin. Test
points on the anomalous side of the boundary receive a negative
decision-function value. The model does not consume fraud labels at training time, which
makes it a natural reference point for quantifying the value added by
supervised labels under extreme imbalance.

\paragraph{FT-Transformer.}
The FT-Transformer (Feature Tokenizer + Transformer)~\cite{gorishniy2021revisiting}
adapts the transformer architecture~\cite{vaswani2017attention} to tabular
inputs. Each input feature --- numerical or categorical --- is first
embedded as a $d$-dimensional token through a feature-specific tokeniser
(numerical features are passed through a learned linear-plus-bias map; each
categorical level has its own embedding vector). A learnable \texttt{[CLS]}
token is prepended to the sequence of feature tokens, and the resulting
sequence of length $T{=}n_{\text{features}}{+}1$ is processed by a stack of
transformer blocks. Each block applies multi-head scaled dot-product
self-attention,
$\mathrm{Attn}(Q,K,V)=\mathrm{softmax}(QK^{\top}/\!\sqrt{d_k})V$, followed by
a position-wise feed-forward network and residual connections with layer
normalisation. The final embedding of the \texttt{[CLS]} token is passed to a
linear classification head. Self-attention provides every feature with a
soft, learned weighting over every other feature, which in principle allows
arbitrary feature interactions to be captured.

% ──────────────────────────────────────────────────────────────────────────
\section{Class-Imbalance Handling Strategies}
\label{sec:tb:imbalance}
% ──────────────────────────────────────────────────────────────────────────

Under extreme class imbalance, the minority class contributes relatively few
examples to the empirical training distribution, and minimising an unmodified
loss can favour the majority class. The
literature on imbalance handling~\cite{hegarcia2009learning, krawczyk2016learning}
groups remedies into three families: \emph{data-level} interventions that
modify the empirical training distribution, \emph{algorithm-level}
interventions that modify the loss, and the unmodified baseline.

\paragraph{Random under- and over-sampling.}
\emph{Random Under-Sampling} (RUS) removes majority examples uniformly at
random until a target ratio is reached. \emph{Random Over-Sampling} (ROS)
duplicates minority examples uniformly at random until the same ratio is
reached. RUS discards information and is therefore most useful when the
majority class is large enough that the resulting training set still provides
a meaningful empirical estimate; ROS does not discard information but, by
duplicating exact examples, can amplify overfitting to the specific minority
realisations present in the training set.

\paragraph{SMOTE and its variants.}
The Synthetic Minority Oversampling Technique~\cite{chawla2002smote}
generates synthetic minority examples by interpolation in feature space: for
each minority point $\mathbf{x}_i$, one of its $k$ nearest minority neighbours
$\mathbf{x}_{j}$ is selected at random and a new example
$\mathbf{x}_i + \lambda(\mathbf{x}_j - \mathbf{x}_i)$ is created with
$\lambda \sim \mathcal{U}(0,1)$. Two cleaning-augmented variants
\cite{batista2004balancing} post-process the SMOTE output to reduce noise
near the decision boundary: \emph{SMOTE+Tomek Links} removes any pair of
nearest neighbours from opposite classes (a Tomek link), and \emph{SMOTEENN}
applies the Edited Nearest Neighbour rule, removing examples whose class
disagrees with the majority of their $k$ nearest neighbours. How many examples
the cleaning step removes depends on the geometry of the resampled training
data. Equal downstream metrics for a cleaning variant and plain SMOTE do not,
by themselves, reveal whether the resampled training sets were identical.

\paragraph{Cost-sensitive learning.}
The cost-sensitive (or \emph{class weighting}) family modifies the training
objective rather than the training data. A common balanced weighting rule
assigns weights inversely proportional to class frequency, increasing the
relative contribution of minority-class errors. The
empirical training distribution is left unchanged, which avoids the
distributional distortions introduced by oversampling, but the model is
encouraged to allocate capacity to the minority class through the reweighted
gradient.

% ──────────────────────────────────────────────────────────────────────────
\section{Evaluation Metrics}
\label{sec:tb:metrics}
% ──────────────────────────────────────────────────────────────────────────

Under extreme class imbalance, the choice of evaluation metric is non-trivial:
several widely used scalars are dominated by the majority class and become
nearly insensitive to the minority. This section defines the metric family
used throughout this thesis.

\paragraph{Confusion-matrix primitives.}
For a binary classifier and a fixed decision rule, the confusion matrix has
four counts: true positives (TP), false positives (FP), true negatives (TN),
and false negatives (FN). Precision and recall are defined as
\[
\mathrm{Precision} = \frac{TP}{TP + FP}, \qquad
\mathrm{Recall} = \frac{TP}{TP + FN}.
\]
Neither expression contains the true-negative count. Precision quantifies the
fraction of alerts that are fraudulent, whereas recall quantifies the fraction
of fraudulent cases detected. They therefore expose positive-class behaviour
that overall accuracy can obscure under extreme imbalance.

\paragraph{The $F_{\beta}$-score.}
The $F_{\beta}$-score is the weighted harmonic mean of precision and recall,
\[
F_{\beta} = (1+\beta^{2})\,\frac{\mathrm{Precision}\cdot\mathrm{Recall}}
                                 {\beta^{2}\,\mathrm{Precision} + \mathrm{Recall}},
\]
with $\beta$ controlling the relative weight assigned to recall: $\beta{=}1$
recovers the symmetric $F_1$, while $\beta{=}2$ ($F_2$) assigns recall four
times the weight of precision through the $\beta^2$ term. In fraud
detection $F_2$ provides a recall-sensitive summary when missed fraud is given
greater importance than a false alert. The value $\beta{=}2$ is a reporting
choice rather than a direct monetary cost model.

\paragraph{PR-AUC and ROC-AUC.}
The precision--recall (PR) curve traces the trade-off between precision and
recall as the decision threshold sweeps the full range of model scores.
This thesis reports average precision under the conventional label PR-AUC as a
threshold-independent summary. The receiver operating characteristic (ROC)
curve traces the true-positive rate against the false-positive rate over the
same threshold sweep, and the area under it is ROC-AUC. ROC-AUC is invariant to
prevalence when the class-conditional score distributions are fixed. Under
extreme imbalance, however, a small false-positive rate can still correspond
to many false alerts, and the ROC curve does not expose precision directly.
PR-AUC, by contrast, has a random-ranking baseline equal to the positive
prevalence~$\pi$ and summarises precision--recall performance for the positive
class; it is therefore the primary ranking metric used in this thesis.

\paragraph{Calibration.}
A probabilistic classifier is well calibrated if its output scores match
empirical conditional positive frequencies: among examples assigned a score
of, say, $0.05$, the fraction that are fraudulent should be approximately $0.05$.
The Brier score $\frac{1}{N}\sum_i(\hat{p}_i - y_i)^2$ is the mean squared
error of the predicted probability against the binary label and is a standard
proper scoring rule for probabilistic accuracy; lower is better. Because it
reflects both calibration and resolution, it is not a pure calibration measure.
Reliability curves assess calibration more directly by plotting empirical
positive frequency against predicted score within score bins.

% ──────────────────────────────────────────────────────────────────────────
\section{Decision-Threshold Selection}
\label{sec:tb:thresholds}
% ──────────────────────────────────────────────────────────────────────────

A classifier produces a score $s(\mathbf{x})$ for each input; for probabilistic
models this score is commonly in $[0,1]$, whereas an anomaly detector may produce
an uncalibrated decision value. Converting a score into a binary decision requires a
\emph{threshold rule} that maps scores to predictions, the canonical form
being $\hat{y} = \mathbb{1}\{s(\mathbf{x}) \geq \tau\}$. The choice of $\tau$
determines where on the precision--recall curve the system operates, and
under extreme imbalance different defensible threshold rules can materially
alter threshold-dependent metrics. Several classes of threshold rule appear
in the literature.

\paragraph{Fixed default ($\tau{=}0.5$).}
The most common choice in software defaults is $\tau{=}0.5$. For a calibrated
posterior probability, this threshold corresponds to a decision rule with
equal false-positive and false-negative costs; it does not imply equal class
prevalence. Under asymmetric operational costs or model-dependent score
calibration, $0.5$ need not be an appropriate operating point. Validation data
should therefore be used to determine whether the default achieves the
required precision--recall trade-off.

\paragraph{Validation-optimal $F_{\beta}$ thresholds.}
A second family selects $\tau$ by maximising a chosen $F_{\beta}$ score on
held-out validation data, $\tau^{*} = \arg\max_{\tau} F_{\beta}(\tau)$.
Setting $\beta{=}1$ gives precision and recall equal weight; $\beta{=}2$
assigns greater weight to recall and is suitable when missed positives receive
greater emphasis. The selected value does not, by itself, encode a monetary
cost ratio.

\paragraph{Precision-constrained thresholds.}
A third family selects $\tau$ to enforce a lower bound on validation precision,
$\tau^{*} = \min\{\tau : \mathrm{Precision}(\tau) \geq P_{\min}\}$. The
resulting operating point targets a minimum positive predictive value and
therefore influences the fraction of false alerts. It does not directly fix the
number of alerts, and the realised precision may differ on future data.

\paragraph{Alert-budget thresholds.}
A fourth family selects $\tau$ as the empirical quantile of the validation
score distribution that flags exactly the top-$k$ scoring instances, which
amounts to fixing the operational alert volume independently of any
performance criterion. This rule represents settings in which investigation
throughput is constrained by a fixed review capacity rather than by a target
precision.

To remain leakage-free, any of these rules must be computed exclusively on
data selected without access to the test set. The model-specific validation procedures
used to select $\tau$ and apply it to the held-out test set are described in
Section~\ref{sec:threshold_study} (p.~\pageref{sec:threshold_study}) of the Methodology chapter.

% ──────────────────────────────────────────────────────────────────────────
\section{Hyperparameter Optimisation}
\label{sec:tb:hpo}
% ──────────────────────────────────────────────────────────────────────────

Hyperparameter optimisation (HPO) selects values for the parameters that
control a learning algorithm but are not learned from data --- regularisation
strength, tree depth, number of attention heads, learning rate, and so on.
Brute-force grid search becomes expensive in high-dimensional configuration
spaces, while random search does not use information from earlier trials.
Modern HPO frameworks can instead adopt
\emph{sequential model-based optimisation} (SMBO), in which a probabilistic
surrogate model of the objective is updated after every completed trial and
the next configuration is selected to maximise an acquisition criterion such
as expected improvement.

\paragraph{Tree-structured Parzen Estimator.}
The Tree-structured Parzen Estimator (TPE), used as the default sampler in
Optuna~\cite{akiba2019optuna}, is one realisation of SMBO. Rather than
modelling the objective $p(y\,|\,\mathbf{x})$ directly, TPE models the two
conditional distributions $p(\mathbf{x}\,|\,y < y^{*})$ and
$p(\mathbf{x}\,|\,y \geq y^{*})$ of "good" and "bad" configurations (split at
a quantile $y^{*}$ of the observed objective values) and proposes the
configuration that maximises the ratio of these two densities. This heuristic
concentrates exploration in regions of the configuration space where the
surrogate estimates that good outcomes are more likely than bad ones.

\paragraph{Pruning of unpromising trials.}
Many HPO objectives produce intermediate values during a single trial --- a
cross-validation fold's metric, an epoch's validation loss --- which a
\emph{pruner} can compare against the trajectory of completed trials. A
median pruner, for example, terminates a trial whose intermediate metric
falls below the median of the completed trials' values at the same step,
allowing the optimiser to spend less computation on unpromising regions.
Pruning reduces the cost of the search but can disadvantage configurations
whose intermediate performance improves slowly.

% ──────────────────────────────────────────────────────────────────────────
\section{Feature Attribution with SHAP}
\label{sec:tb:shap}
% ──────────────────────────────────────────────────────────────────────────

SHAP (\emph{SHapley Additive exPlanations})~\cite{lundberg2017unified}
attributes a model's prediction $f(\mathbf{x})$ on a specific instance to its
input features by treating the features as cooperating players in a
coalitional game whose payoff is the prediction itself. The SHAP value
$\phi_{j}$ assigned to feature $j$ is the cooperative-game-theoretic Shapley
value of that feature: the average marginal contribution of feature $j$ to
the prediction across all possible orderings in which features can be
revealed to the model,
\[
\phi_{j} = \sum_{S \subseteq F\setminus\{j\}} \frac{|S|!\,(|F|-|S|-1)!}{|F|!}
           \Bigl[ f_{S\cup\{j\}}(\mathbf{x}_{S\cup\{j\}}) - f_{S}(\mathbf{x}_{S}) \Bigr].
\]
Within the additive explanation framework, the formulation is characterised by
three properties: \emph{local accuracy} (the attributions sum to the prediction
relative to a baseline), \emph{missingness} (a missing feature receives zero
attribution), and \emph{consistency} (an attribution cannot decrease when that
feature's marginal contribution increases across coalitions). These properties
provide a common vocabulary for model explanation, but comparisons across
models still depend on compatible output scales, background data, feature
grouping, and explainer assumptions.

\paragraph{Architecture-specific explainers.}
Computing the Shapley value exactly requires evaluating the model over all
$2^{|F|}$ feature subsets, which is intractable for non-trivial feature counts.
Architecture-specific explainers exploit model structure to make the
computation efficient. \emph{LinearExplainer} uses the linear model form and a
specified background distribution. \emph{TreeSHAP}~\cite{lundberg2020local}
computes exact values for its tree-specific SHAP formulation in polynomial
time. \emph{GradientExplainer} approximates expected gradients for
differentiable models using a background sample; its output depends on both
the approximation and the selected background data.

\paragraph{Comparing attributions across models.}
Once per-feature SHAP values are available, a model's global feature
importance is summarised as the mean absolute SHAP value across a held-out
sample. Two models can then be compared on the agreement between their
top-$k$ feature sets through the \emph{Jaccard similarity}
$|A \cap B|/|A \cup B|$, which equals one when the two sets coincide and
zero when they are disjoint. In this thesis, Jaccard agreement is used only to
compare membership of the top-$k$ global feature sets. It does not compare rank
order, attribution magnitude, local explanations, or whether different models
have learned the same decision process.

% ──────────────────────────────────────────────────────────────────────────
\section{Chapter Summary}
\label{sec:tb:summary}
% ──────────────────────────────────────────────────────────────────────────

This chapter has defined the model families, imbalance interventions,
evaluation measures, threshold rules, optimisation procedure, and SHAP
quantities used in the empirical study. Chapter~\ref{sec:chapterMethodology}
now translates these general concepts into the project-specific datasets,
preprocessing paths, validation procedures, and analysis plan.
